% ============================================================================
% Collaboration Networks in a Tabula Rasa Legislature — presentación 23-jul-2026
% v2 (correcciones del autor 2026-07-22): RQ1 reorientada como comparación
% CLogit vs ERGM bipartito; sección de datos ampliada; listas/Rice y
% "credenciales" fuera; RHEM a backup; títulos de lámina más arriba.
% Estilo replicado de presentation-cpin2.tex (JORNADA REDES).
% Compilar: pdflatex presentation-2026-07-23.tex  (correr dos veces)
% ============================================================================
\documentclass[aspectratio=169,11pt]{beamer}

\usepackage[utf8]{inputenc}
\usepackage[T1]{fontenc}
\usepackage{amsmath, amssymb}
\usepackage{graphicx}
\usepackage{booktabs}
\usepackage{array}
\usepackage{xcolor}
\usepackage{qrcode}

% ---------- Fuentes ----------
\usefonttheme{professionalfonts}
\usepackage{helvet}
\renewcommand{\familydefault}{\sfdefault}

% ---------- Colores (deck JORNADA) ----------
\definecolor{ndred}{RGB}{222,32,38}
\definecolor{ndreddark}{RGB}{176,22,28}
\definecolor{ndgreen}{RGB}{0,118,54}
\definecolor{ndblue}{RGB}{30,90,200}
\definecolor{ndgray}{RGB}{90,90,90}

% ---------- Sin chrome de Beamer ----------
\setbeamertemplate{navigation symbols}{}
\setbeamercolor{normal text}{fg=black,bg=white}
\setbeamercolor{structure}{fg=ndreddark}
\setbeamertemplate{itemize item}{\color{ndreddark}\textbullet}
\setbeamertemplate{itemize subitem}{\color{ndgray}\textendash}

% ---------- Título de lámina: centrado + regla roja, PEGADO ARRIBA ----------
\setbeamercolor{frametitle}{fg=black,bg=white}
\setbeamerfont{frametitle}{size=\large, series=\mdseries}
\setbeamertemplate{frametitle}{%
  \vspace{0.12cm}%
  \begin{center}\insertframetitle\end{center}%
  \vspace{-0.20cm}%
  {\color{ndred}\hrule height 1.1pt}%
}

% ---------- Pie mínimo: número de página ----------
\setbeamertemplate{footline}{%
  \hfill{\color{ndgray}\scriptsize\insertframenumber}\hspace{0.5cm}\vspace{0.2cm}}

% ---------- Macros ----------
\newcommand{\emphr}[1]{{\color{ndreddark}\textbf{#1}}}
\newcommand{\nota}[1]{{\footnotesize\color{ndgray} #1}}
\newcommand{\sectionslide}[2]{%
  \begin{frame}[plain]
    \vfill
    \begin{center}
      {\footnotesize\color{ndgray} #1}\\[6pt]
      {\Large #2}\\[2pt]
      {\color{ndred}\rule{0.30\textwidth}{1.1pt}}
    \end{center}
    \vfill
  \end{frame}}
\newcommand{\fig}[2]{\begin{center}\includegraphics[width=#1\textwidth]{../results/figures/#2}\end{center}}

% ============================================================================
\title{Collaboration Networks in a Tabula Rasa Legislature}
\author{Aníbal Olivera M. \& Jorge Fábrega}
\date{}

\begin{document}

% ---------------------------------------------------------------- PORTADA
\begin{frame}[plain]
  \vspace{0.6cm}
  \begin{center}
    {\footnotesize\color{ndgray} Centro de Investigación en Complejidad Social $\cdot$ Universidad del Desarrollo $\cdot$ Chile}
  \end{center}
  \vspace{1.0cm}
  \begin{center}
    {\LARGE \textbf{Collaboration Networks in a Tabula Rasa Legislature}}\\[0.5cm]
    {\large \color{ndreddark} Formación de lazos, conducta y éxito legislativo\\en la Convención Constitucional chilena (2021--2022)}
  \end{center}
  \vspace{1.1cm}
  \begin{center}
    {\normalsize Aníbal Olivera M. \quad$\cdot$\quad Jorge Fábrega}\\[2pt]
    {\footnotesize PhD(c) in Social Complexity Sciences, CICS--UDD}
  \end{center}
  \vfill
  \begin{center}
    {\scriptsize\color{ndgray} 23 de julio de 2026}
  \end{center}
\end{frame}

% ---------------------------------------------------------------- ÍNDICE
\begin{frame}{Índice}
\begin{itemize}
  \item[{\color{ndred}\textbf{1.}}] La Convención: el cuasi-experimento y los datos
  \item[{\color{ndred}\textbf{2.}}] RQ1 — Formación: ¿podíamos predecir la red?
  \begin{itemize}
    \item[{\color{ndred}2.1}] {\color{black}Herramienta 1: el logit condicional (elección dado el menú)}
    \item[{\color{ndred}2.2}] {\color{black}Herramienta 2: el ERGM bipartito (la estructura completa)}
    \item[{\color{ndred}2.3}] {\color{black}Dónde coinciden las lentes, y dónde no}
  \end{itemize}
  \item[{\color{ndred}\textbf{3.}}] ¿Qué compra el capital relacional?
  \begin{itemize}
    \item[{\color{ndred}3.1}] {\color{black}RQ2a — ¿La red mueve las posiciones?}
    \item[{\color{ndred}3.2}] {\color{black}RQ2b — ¿La red mueve la conducta?}
    \item[{\color{ndred}3.3}] {\color{black}RQ3 — ¿Qué hace que un artículo sobreviva?}
  \end{itemize}
  \item[{\color{ndred}\textbf{4.}}] Cierre: la aritmética del 103
\end{itemize}
\end{frame}

% ---------------------------------------------------------------- APERTURA
\begin{frame}{Un cuasi-experimento natural}
\begin{itemize}
  \item En una legislatura ordinaria, las redes de colaboración son el \emphr{sedimento} de décadas: carreras, comités, favores, disciplinas de partido. Imposible ver cómo \emph{nacen}.
  \vspace{0.35cm}
  \item La Convención Constitucional chilena (2021--22) ofrece lo que casi nunca se observa: \emphr{colaboración política organizándose desde cero}.
  \vspace{0.35cm}
  \item Mayoría de novatos, órgano que se disuelve al terminar, reglas escritas por sus propios miembros — y \emphr{registro documental completo} de cada firma, enmienda y voto.
  \vspace{0.35cm}
  \item Lo que emerja, emergió de lo que la gente \emphr{traía puesto}. Ese es el experimento.
\end{itemize}
\end{frame}

% ---------------------------------------------------------------- RQs
\begin{frame}{Las preguntas de investigación}
\begin{itemize}
  \item El registro documental completo permite muchas preguntas (formación, agenda, género, especialización, \dots). Hoy: \emphr{las de redes}.
  \vspace{0.3cm}
  \item \emphr{RQ1 — Formación}: ¿podíamos \emph{predecir} la red de colaboración solo con lo que se sabía de cada convencional \emph{antes} del día uno?
  \vspace{0.3cm}
  \item ¿Qué compra el capital relacional?
  \begin{itemize}
    \item \emphr{RQ2a — Posiciones}: ¿la red mueve las \emph{posiciones} de los convencionales?
    \item \emphr{RQ2b — Conducta}: ¿la red coordina su \emph{conducta} de voto?
    \item \emphr{RQ3 — Textos}: ¿qué hace que un artículo \emph{sobreviva} hasta el borrador?
  \end{itemize}
\end{itemize}
\end{frame}

% ================================================================ PARTE 0
\sectionslide{Parte 0}{La Convención}

\begin{frame}{La Convención: Tabula Rasa}
\begin{itemize}
  \item 155 electos (154 en ejercicio), julio 2021 -- julio 2022; mayoría inédita de independientes y novatos; 17 escaños reservados para pueblos originarios; paridad de género.
  \vspace{0.25cm}
  \item El órgano \emphr{se disolvía al entregar el texto}: sin sombra electoral del futuro.
  \vspace{0.25cm}
  \item Dos reglas, escritas por los propios convencionales, estructuran el estudio:
  \begin{itemize}
    \item Cada iniciativa constituyente requería \emphr{8--16 patrocinantes}: firmar = formar una coalición visible, fechada y acotada.
    \item Cada norma requería \emphr{2/3 del Pleno (103 de 154)}; en comisión bastaba mayoría simple.
  \end{itemize}
  \vspace{0.25cm}
  \item[] \nota{Nacer era barato; sobrevivir, caro.}
\end{itemize}
\end{frame}

\begin{frame}{El proceso en el tiempo}
\fig{0.97}{cc_timeline.pdf}
\end{frame}

\begin{frame}{Los datos: el proceso documental completo}
\begin{itemize}
  \item \emphr{Actores}: los 154 convencionales, con perfil curado por doble fuente (BCN + Wikipedia): profesión (39\% abogados/as), experiencia institucional previa, educación (0--3), edad, género, distrito/pueblo, lista, comisión.
  \item \emphr{Iniciativas}: 995 ingresadas a la plataforma oficial (nov-21 a feb-22); \emphr{947} con 2--16 firmantes-persona forman el set de análisis; \emphr{100\% fechadas}.
  \item \emphr{Trazabilidad de textos}: 1.565 artículos génesis con desenlace conocido frente al borrador (supervivencia 20\%), con historial fechado de indicaciones.
  \item \emphr{Indicaciones}: 2.326 actos de enmienda con autoría; 1.004 eventos multi-firmantes.
  \item \emphr{Votaciones}: 4.707 votaciones nominales del Pleno.
\end{itemize}
\vspace{0.15cm}
\nota{Pipeline reproducible; los datos viven como snapshot versionado. La reconstrucción de este registro (armonización de 75.616 menciones de nombres, cruce plataforma-informes) es parte de la contribución.}
\end{frame}

\begin{frame}{Las siete comisiones temáticas}
\begin{center}\scriptsize
\begin{tabular}{clccccccccc}
\toprule
 & Nombre (corto) & Miembros & \% abog. & \% exper. & Edad & Grado & Iniciativas & Ondas & Indicaciones & Multifirm. \\
\midrule
C1 & Sistema Político & 25 & 60 & 36 & 42.2 & 1.44 & 89 & 4 & 365 & 266 \\
C2 & Principios Const. & 18 & 22 & 22 & 45.8 & 1.28 & 114 & 3 & 53 & 1 \\
C3 & Forma de Estado & 30 & 37 & 40 & 45.1 & 1.43 & 51 & 6 & 241 & 212 \\
C4 & Derechos Fund. & 30 & 30 & 13 & 45.9 & 1.03 & 283 & 5 & 395 & 154 \\
C5 & Medio Ambiente & 19 & 21 & 16 & 44.7 & 1.26 & 133 & 5 & 559 & 0 \\
C6 & Sist. de Justicia & 17 & 88 & 18 & 42.6 & 1.65 & 93 & 6 & 381 & 317 \\
C7 & Conocimientos & 15 & 13 & 0 & 51.7 & 1.33 & 64 & 8 & 332 & 54 \\
\bottomrule
\end{tabular}
\end{center}
\vspace{0.2cm}
\begin{itemize}
  \item Composición y productividad \emphr{fuertemente heterogéneas} — los modelos las explotan y controlan.
  \item \nota{Iniciativas = con 2--16 firmantes asignadas a la comisión (otras 120 sin comisión asignada). Multifirm. = indicaciones con $\geq 2$ firmantes, las únicas que agregan lazos. El cero de C5 es limitación de registro.}
\end{itemize}
\end{frame}

\begin{frame}{Red de colaboración: 1}
\fig{0.84}{bipartite_initiative.pdf}
\vspace{0.05cm}
\nota{\emphr{Esta es la unidad de análisis real}: convencional $\times$ documento — una firma, un lazo. No partimos de una red actor--actor: la proyección persona-persona es una \emph{decisión de modelamiento} (una iniciativa de 16 firmantes fabricaría 120 pares de una sola vez).}
\end{frame}

\begin{frame}{Red de colaboración: 2}
\fig{0.95}{indication_events_timeline.pdf}
\vspace{0.05cm}
\nota{El segundo reloj: la enmienda. Cada barra es un día de informe de comisión; sólido = eventos multi-firmantes (los que agregan lazos a la red dinámica); claro = total fechado.}
\end{frame}

\begin{frame}{Cómo medimos la ideología}
\begin{itemize}
  \item \emphr{Pre-red}: W-NOMINATE 2D con \emph{solo el primer mes} de votaciones del Pleno (147, jul--ago 2021), replicando a Fábrega (2022, \emph{RCP}).
  \begin{itemize}
    \item Esa ventana es anterior a las comisiones (oct), a las iniciativas (nov) y a la regla de 2/3 operativa (feb) $\Rightarrow$ $\theta_1, \theta_2$ son covariables \emphr{exógenas a la red que estudiamos}.
    \item De aquí sale el \emphr{pívot de 2/3}: ordenados los 154, el convencional 103 marca $\theta_{1,(103)} = -0.15$.
  \end{itemize}
  \vspace{0.3cm}
  \item \emphr{Voto revelado dinámico}: dynIRT sobre las 4.707 votaciones $\to$ trayectoria $\theta_{i,t}$ en 91 períodos (insumo de la parte de influencia; se lee como conducta de voto, no como ideología latente pura).
\end{itemize}
\end{frame}

\begin{frame}{Las posiciones en el tiempo, por comisión}
\fig{0.88}{positions_dynamics_by_commission.pdf}
\end{frame}

% ================================================================ ACTO I
\sectionslide{RQ1 — Formación}{¿Podíamos predecir la red de colaboración?}

\begin{frame}{La pregunta, y las dos herramientas}
\begin{itemize}
  \item Si la Convención es un cuasi-experimento, la pregunta natural es: ¿cuánto de la red de colaboración se \emphr{predice con lo que ya se sabía} de cada convencional antes del día uno?
  \vspace{0.25cm}
  \item La respuesta \emphr{depende de la herramienta} — y esa es parte de la lección:
  \begin{itemize}
    \item \emphr{Logit condicional}: cada iniciativa como un \emph{menú}; compara firmantes contra no-firmantes \emph{dentro} del mismo menú. Condiciona en que la iniciativa existe y en su tamaño.
    \item \emphr{ERGM bipartito}: modela la \emph{estructura completa} convencional $\times$ documento — la unidad real del dato — sin proyectar ni condicionar.
  \end{itemize}
  \vspace{0.25cm}
  \item Veremos qué covariables sobreviven a ambas lentes, cuáles cambian de cara — y qué se aprende de usar la herramienta que aprovecha toda la riqueza del dato.
\end{itemize}
\end{frame}

\begin{frame}{Herramienta 1: el logit condicional (elección dado el menú)}
\begin{itemize}
  \item Utilidad latente de firmar; el atractivo propio del texto ($\alpha_a$: tema, redactor, momento) \emphr{se cancela} al comparar dentro del menú:
\end{itemize}
$$U_{ia} = \beta^\top x_{ia} + \alpha_a + \varepsilon_{ia}
\qquad\Rightarrow\qquad
P(S_a \mid |S_a|) = \frac{\exp(\sum_{i \in S_a}\beta^\top x_{ia})}{\sum_{|R|=|S_a|}\exp(\sum_{i \in R}\beta^\top x_{ia})}$$
\begin{itemize}
  \item 947 menús $\times$ 154 candidatos = 145.838 decisiones; covariables leave-one-out respecto de la coalición; EE cluster por convencional.
  \item \nota{Fortaleza: inferencia exacta y control total de composición del menú. Costo: solo ve la elección \emph{dado} que el evento existe — no modela la estructura.}
\end{itemize}
\end{frame}

\begin{frame}{Qué predice la firma según el logit condicional}
\begin{center}\scriptsize
\begin{tabular}{lcccc}
\toprule
Variable & Coef. & OR & EE & $p$ \\
\midrule
\multicolumn{5}{l}{\emphr{--- Controles: estructura y posiciones formadas en la Convención}}\\
Misma comisión que el texto & $+0.93$ & 2.5 & 0.05 & $<10^{-72}$ \\
Distancia en $\theta_1$ a la coalición & $-3.64$ & 0.03 & 0.24 & $<10^{-51}$ \\
Distancia en $\theta_2$ a la coalición & $-1.38$ & 0.25 & 0.14 & $<10^{-20}$ \\
\multicolumn{5}{l}{\emphr{--- Listas electorales (pre-Convención)}}\\
Misma lista: Escaños Reservados PPOO & $+1.43$ & 4.2 & 0.21 & $<10^{-10}$ \\
Misma lista: Otras listas locales & $+1.33$ & 3.8 & 0.16 & $<10^{-16}$ \\
Misma lista: Lista del Pueblo & $+1.07$ & 2.9 & 0.16 & $<10^{-10}$ \\
Misma lista: Lista del Apruebo & $+1.05$ & 2.9 & 0.19 & $<10^{-7}$ \\
Misma lista: Vamos por Chile & $+0.94$ & 2.6 & 0.21 & $<10^{-4}$ \\
Misma lista: Apruebo Dignidad & $+0.93$ & 2.5 & 0.14 & $<10^{-9}$ \\
\multicolumn{5}{l}{\emphr{--- Territorio (pre-Convención)}}\\
Afinidad de distrito/pueblo & $+2.52$ & 12.4 & 0.56 & $<10^{-5}$ \\
\multicolumn{5}{l}{\emphr{--- Perfil pre-Convención}}\\
Afinidad de abogados & $+0.12$ & 1.1 & 0.18 & $0.48$ \\
Afinidad de experiencia previa & $+0.52$ & 1.7 & 0.23 & $0.024$ \\
Distancia de grado académico & $-0.10$ & 0.9 & 0.08 & $0.21$ \\
Distancia de edad (décadas) & $-0.07$ & 0.9 & 0.06 & $0.23$ \\
Afinidad de género & $+0.28$ & 1.3 & 0.16 & $0.093$ \\
\midrule
\multicolumn{5}{l}{AIC = 87.317; \quad pseudo-$R^2$ (McFadden) = 0.209}\\
\bottomrule
\end{tabular}
\end{center}
\end{frame}

\begin{frame}{La heterogeneidad por bloque del elector (mismo modelo, submuestras)}
\begin{center}\footnotesize
\begin{tabular}{lccccc}
\toprule
 & Derecha & Centro-izq & Izquierda & PPOO & Otras \\
 & (VC, 37) & (LA+INN, 28) & (AD+LdP, 51) & (17) & (21) \\
\midrule
Afinidad distrito/pueblo & $+7.9^{***}$ & $+5.1^{***}$ & $+0.1$ & $-1.4^{+}$ & $+2.4^{+}$ \\
Afinidad de género & $+0.6^{*}$ & $+1.9^{***}$ & $+0.3$ & $+0.0$ & $-0.0$ \\
Afinidad de abogados & $+0.0$ & $+0.5$ & $-0.0$ & $+0.9$ & $+1.2^{*}$ \\
Distancia de edad & $+0.1$ & $-0.1$ & $-0.2^{**}$ & $-0.1$ & $+0.1$ \\
Distancia en $\theta_1$ & $-5.1^{***}$ & $-2.8^{**}$ & $-0.8$ & $+0.1$ & $-6.1^{***}$ \\
\bottomrule
\end{tabular}
\end{center}
\vspace{0.25cm}
\begin{itemize}
  \item El territorio es de \emphr{derecha y centro}; la izquierda organizada coordina por lista e ideología, no por geografía; el género vive en la centro-izquierda.
  \item La pendiente ideológica: la derecha casi no cruza distancias; la izquierda las cruza todas.
\end{itemize}
\end{frame}

\begin{frame}{Herramienta 2: el ERGM bipartito (la estructura completa)}
\begin{itemize}
  \item Modelo de probabilidad sobre la red convencional $\times$ documento entera:
\end{itemize}
$$\Pr(Y = y) = \frac{\exp\big(\theta^\top s(y)\big)}{\kappa(\theta)}$$
\nota{$s(y)$: firmas de miembros, pares de co-firmantes que comparten atributo, \dots}
\begin{itemize}
  \item Cada término pregunta: ¿este patrón aparece \emphr{más que en redes aleatorias comparables}?
  \item \emphr{Siete modelos, uno por comisión}: la comisión es el mayor confundidor de composición (el modelo agregado con intercepto único invierte los signos — Simpson verificado).
  \item \nota{El denominador $\kappa(\theta)$ suma sobre \emph{todas} las redes posibles ($2^{154 \times n}$) — incalculable. MCMC lo \emph{aproxima simulando} redes (aquí resultó infactible); \emphr{MPLE lo esquiva}: estima cada lazo condicional al resto (una logística sobre estadísticas de cambio) — mismos puntos en segundos; los EE honestos vienen de \emphr{bootstrap de iniciativas}.}
\end{itemize}
\end{frame}

\begin{frame}{La misma pregunta con la estructura completa}
\begin{center}\tiny
\begin{tabular}{lccccccc}
\toprule
Término & C1 & C2 & C3 & C4 & C5 & C6 & C7 \\
\midrule
\multicolumn{8}{l}{\emphr{--- Controles: base, pertenencia e ideología}}\\
Edges (base) & $-3.14^{***}$ & $-3.90^{***}$ & $-3.70^{***}$ & $-3.60^{***}$ & $-3.78^{***}$ & $-3.38^{***}$ & $-3.79^{***}$ \\
Miembro de la comisión & $+0.73^{***}$ & $+1.24^{***}$ & $+1.77^{***}$ & $+0.57^{***}$ & $+1.56^{***}$ & $+1.31^{***}$ & $+2.29^{***}$ \\
Mismo quintil $\theta_1$ & $+0.16^{***}$ & $+0.16^{***}$ & $+0.12^{***}$ & $+0.14^{***}$ & $+0.16^{***}$ & $+0.13^{***}$ & $+0.16^{***}$ \\
\multicolumn{8}{l}{\emphr{--- Listas electorales}}\\
Misma lista & $+0.21^{***}$ & $+0.18^{***}$ & $+0.20^{***}$ & $+0.18^{***}$ & $+0.18^{***}$ & $+0.21^{***}$ & $+0.22^{***}$ \\
\multicolumn{8}{l}{\emphr{--- Territorio}}\\
Mismo distrito/pueblo & $+0.23^{***}$ & $+0.19^{***}$ & $+0.36^{***}$ & $-0.01$ & $+0.07$ & $-0.02$ & $+0.30^{***}$ \\
\multicolumn{8}{l}{\emphr{--- Perfil pre-Convención}}\\
Ambos abogados & $-0.08^{***}$ & $-0.02^{+}$ & $-0.08^{***}$ & $-0.01^{+}$ & $-0.03^{**}$ & $-0.02$ & $-0.01$ \\
Ambos con experiencia & $+0.04^{**}$ & $+0.06^{***}$ & $+0.05^{**}$ & $+0.06^{***}$ & $+0.06^{***}$ & $+0.05^{***}$ & $+0.02$ \\
Mismo género & $-0.03^{+}$ & $-0.00$ & $-0.05^{*}$ & $-0.01$ & $-0.06^{***}$ & $-0.08^{***}$ & $-0.04^{+}$ \\
Mismo nivel educativo & $-0.10^{***}$ & $-0.08^{***}$ & $-0.03^{+}$ & $-0.09^{***}$ & $+0.02^{*}$ & $-0.05^{***}$ & $+0.00$ \\
Mismo quintil de edad & $-0.06^{*}$ & $-0.04^{+}$ & $-0.01$ & $-0.05^{***}$ & $-0.08^{***}$ & $-0.05^{*}$ & $-0.09^{**}$ \\
\midrule
\multicolumn{8}{l}{Rangos continuos por documento (\texttt{b2covrange}): $\theta_1 \in [-9.7, -3.3]$, $p<.001$ en las 7 (\emphr{evitan} dispersión ideológica); grado $-$ en 5/7; edad $\approx 0$.}\\
\bottomrule
\end{tabular}
\end{center}
\vspace{0.1cm}
\nota{Homofilia de co-firma medida sin proyección, dentro de cada pool temático.}
\end{frame}

\begin{frame}{Dónde coinciden las lentes, y dónde no}
\begin{center}\footnotesize
\begin{tabular}{lccl}
\toprule
Covariable & CLogit (menú) & ERGM bipartito (estructura) & Lectura \\
\midrule
Misma lista & $+0.9$ a $+1.4^{***}$ & $+0.18$ a $+0.22^{***}$ en las 7 & \emphr{coinciden} \\
Ideología ($\theta_1$) & $-3.6^{***}$ & quintil $+^{***}$; rango $-3$ a $-10^{***}$ & \emphr{coinciden} (más nítida) \\
Distrito/pueblo & $+2.5^{***}$ (global) & fuerte en C1/C2/C3/C7; nulo C4/C5/C6 & el bipartito \emphr{matiza} \\
Experiencia & $+0.5^{*}$ & $+0.02$ a $+0.06$, parejo & coinciden (chico) \\
Abogado & $+0.1$ (n.s.) & levemente \emphr{negativo} (6/7) & difieren de signo \\
Género & $+0.3^{+}$ & \emphr{negativo} casi en las 7 & difieren: \emph{mezcla} \\
Educación & $-0.1$ (n.s.) & match $-$ y rango $-$ & el bipartito \emphr{revela} mezcla moderada \\
Edad & $-0.1$ (n.s.) & quintil $-$; rango $\approx 0$ & mezcla generacional \\
\bottomrule
\end{tabular}
\end{center}
\vspace{0.25cm}
\begin{itemize}
  \item Aprendizaje 1 — \emphr{qué importa}: lista, ideología y (condicionado a comisión) territorio predicen la red; las credenciales no la tejen — y en la estructura completa, las coaliciones \emph{mezclan} género, educación y generación.
  \item Aprendizaje 2 — \emphr{la herramienta importa}: la unidad real (persona $\times$ documento) revela signos y heterogeneidades que el menú promedia.
\end{itemize}
\end{frame}

% ================================================================ QUE COMPRA
\begin{frame}[plain]
  \vfill
  \begin{center}
    {\Large ¿Qué compra el capital relacional?}\\[2pt]
    {\color{ndred}\rule{0.30\textwidth}{1.1pt}}
  \end{center}
  \vfill
\end{frame}

\sectionslide{RQ2 — Personas}{¿La red mueve las posiciones (RQ2a), o la conducta (RQ2b)?}

\begin{frame}{RQ2a: ¿la red mueve las posiciones? El diseño}
\begin{itemize}
  \item Exposición de $i$ en la onda $t$: posición media de sus co-firmantes, ponderada por intensidad:
\end{itemize}
$$E_{i,t} = \frac{\sum_{j \neq i} w_{ij,t}\,\theta_{j,t}}{\sum_{j \neq i} w_{ij,t}}
\qquad\qquad
\Delta\theta_{i,t} = \alpha_i + \beta\,\theta_{i,t-1} + \lambda\,E_{i,t-1} + \varepsilon_{it}$$
\begin{itemize}
  \item Efectos fijos individuales (solo variación \emph{dentro} de cada persona), EE cluster; 4.355 obs.
  \item Resultado: $\hat\lambda = +0.007$ (EE $0.004$), \emphr{nulo}. Robusto a 3 ventanas temporales $\times$ 7 dosis de exposición.
\end{itemize}
\end{frame}

\begin{frame}{Selección, no influencia — en una imagen}
\fig{0.88}{m2_selection_vs_influence_preview.pdf}
\vspace{0.05cm}
\nota{(a) Entre personas: posición y vecindario van pegados ($r=0.95$) — elegimos vecindarios que se nos parecen. (b) Dentro de cada persona, los cambios no se siguen. Ámbar: el 5\% más desalineado — los únicos con espacio para ser arrastrados — cae en (b) sin pendiente ($r=-0.04$).}
\end{frame}

\begin{frame}{La tabla completa: tres ventanas $\times$ siete dosis de exposición}
\begin{center}\tiny
\begin{tabular}{llcccc}
\toprule
Ventana & Definición de exposición & $\hat\lambda$ & EE & $p$ & N \\
\midrule
Completa & Acumulada desde T0 & $+0.007$ & 0.004 & $0.119$ & 4.355 \\
Completa & Solo última onda & $+0.007$ & 0.013 & $0.60$ & 463 \\
Completa & Últimas 2 ondas & $+0.003$ & 0.009 & $0.72$ & 1.504 \\
Completa & Últimas 3 ondas & $-0.002$ & 0.007 & $0.76$ & 2.447 \\
Completa & Decaimiento $\lambda_w = 0.25$ & $+0.008$ & 0.004 & $0.051$ & 4.355 \\
Completa & Decaimiento $\lambda_w = 0.50$ & $+0.008$ & 0.004 & $0.042$ & 4.355 \\
Completa & Decaimiento $\lambda_w = 0.75$ & $+0.007$ & 0.004 & $0.078$ & 4.355 \\
Completa & Falsificación: exposición futura & $+0.010$ & 0.004 & $0.030$ & 3.329 \\
Temprana ($\leq$ 31-mar) & Acumulada & $+0.005$ & 0.007 & $0.50$ & 1.586 \\
Temprana & Últimas 2 ondas & $+0.013$ & 0.014 & $0.35$ & 955 \\
Temprana & Últimas 3 ondas & $+0.002$ & 0.010 & $0.81$ & 1.466 \\
Temprana & Decaimiento $\lambda_w = 0.25$ & $+0.005$ & 0.007 & $0.46$ & 1.586 \\
Temprana & Decaimiento $\lambda_w = 0.50$ & $+0.005$ & 0.007 & $0.48$ & 1.586 \\
Temprana & Decaimiento $\lambda_w = 0.75$ & $+0.005$ & 0.007 & $0.49$ & 1.586 \\
Temprana & Falsificación: exposición futura & $+0.006$ & 0.007 & $0.41$ & 1.587 \\
Tardía ($\geq$ 1-abr) & Acumulada & $+0.011$ & 0.005 & $0.018$ & 2.769 \\
Tardía & Solo última onda & $-0.002$ & 0.018 & $0.92$ & 332 \\
Tardía & Últimas 2 ondas & $-0.006$ & 0.008 & $0.45$ & 549 \\
Tardía & Últimas 3 ondas & $-0.011$ & 0.008 & $0.21$ & 981 \\
Tardía & Decaimiento $\lambda_w = 0.25$ & $+0.009$ & 0.004 & $0.032$ & 2.769 \\
Tardía & Decaimiento $\lambda_w = 0.50$ & $+0.011$ & 0.004 & $0.008$ & 2.769 \\
Tardía & Decaimiento $\lambda_w = 0.75$ & $+0.011$ & 0.004 & $0.011$ & 2.769 \\
Tardía & Falsificación: exposición futura & $+0.010$ & 0.004 & $0.004$ & 1.742 \\
\bottomrule
\end{tabular}
\end{center}
\vspace{0.05cm}
\nota{Ningún $\hat\lambda$ supera $0.011$; los pocos significativos de la ventana tardía vienen acompañados de una falsificación igual de significativa y de la misma magnitud.}
\end{frame}

\begin{frame}{Tres razones para creerle al nulo}
\begin{enumerate}
  \item \emphr{Había espacio para ver influencia.} Podría ser que no hubiera nada que copiar: si mi vecindario ya piensa exactamente como yo, la influencia no deja huella. Lo medimos: la distancia promedio entre cada convencional y su vecindario es $0.59$ (contra $2.25$ si los vecindarios fueran al azar). La selección cerró el 74\% del espacio — pero el $0.59$ restante es espacio real donde un arrastre se habría notado.
  \vspace{0.25cm}
  \item \emphr{Había instrumento para verla.} ¿Y si el efecto existe pero nuestros datos no alcanzan a detectarlo? El efecto mínimo detectable — el $\lambda$ más chico que este diseño habría distinguido de cero — es $0.011$: una influencia que cerrara apenas el 1\% de la distancia por onda. Lo estimado ($0.007$) queda \emph{debajo} incluso de eso.
  \vspace{0.25cm}
  \item \emphr{El termómetro no es el culpable.} $\theta$ se estima desde los votos, con error (grande: 3.6$\times$ el movimiento típico entre períodos). Lo propagamos: re-simulamos las votaciones 50 veces, re-estimamos $\theta$ y el modelo completo en cada réplica. El error estándar de $\lambda$ apenas se mueve ($0.0036 \to 0.0039$), porque el ruido individual se promedia sobre 154 personas y todas las ondas: el nulo no es un artefacto de medición.
\end{enumerate}
\end{frame}

\begin{frame}{RQ2b: la conducta — qué es defección, y qué mide $\phi$}
\begin{itemize}
  \item \emphr{Defección}: en una votación $v$, votar \emph{distinto} de la mayoría del propio bloque electoral. Casi siempre se respeta el libreto: ocurre en el \emphr{7.9\%} de los votos.
  \item La pregunta: cuando alguien se sale del libreto, ¿se sale \emph{solo}, o acompañado de la gente con la que \emph{escribió iniciativas}? Exposición de $i$ en $v$:
\end{itemize}
$$X_{iv} = \frac{\sum_{j \neq i} w_{ij}\, D_{jv}}{\sum_{j \neq i} w_{ij}}
\qquad\qquad
\Pr(D_{iv} = 1) = \Lambda\big(\eta_i + \mu_v + \phi\, X_{iv}\big)$$
\begin{itemize}
  \item $w_{ij}$ = iniciativas co-firmadas; $\eta_i$ absorbe personalidades díscolas; $\mu_v$ absorbe votaciones que rompen a todos.
  \item \emphr{$\phi$} mide: ¿cuánto sube mi propensión (en log-odds) a defeccionar cuando la fracción ponderada de mis co-firmantes que defecciona \emph{en esta misma votación} pasa de 0 a 1?
  \item El riesgo: si una votación parte a mi lista, defeccionamos juntos \emph{sin conocernos}. El antídoto: \emphr{permutación} — barajar los nombres de los defectores dentro de cada bloque $\times$ votación y re-estimar 200 veces.
\end{itemize}
\end{frame}

\begin{frame}{RQ2b: la tabla completa}
\begin{center}\footnotesize
\begin{tabular}{llccc}
\toprule
Modelo & Variable & Coef. & EE & $p$ \\
\midrule
Principal & Exposición a defectores ($\phi$) & $+11.21$ & 0.52 & $<10^{-100}$ \\
Benchmark permutado (200 réplicas) & $\phi$ esperado por mecánica & $6.02$ [p95: $6.08$] & --- & $<0.005$ \\
Período completo (robustez) & $\phi$ & $+11.06$ & 0.47 & $<10^{-121}$ \\
Split por comisión & Exposición co-firmantes de \emph{otra} comisión & $+8.67$ & 0.46 & $<10^{-15}$ \\
 & Exposición co-firmantes de \emph{mi} comisión & $+2.99$ & 0.32 & $<10^{-15}$ \\
 & Tasa de defección de mi comisión & $-3.01$ & 0.45 & $<10^{-10}$ \\
Rezagada: votación anterior (sola) & Exposición en $v-1$ & $+2.98$ & 0.17 & $<10^{-65}$ \\
Horse race & Contemporánea ($\phi$) & $+11.08$ & 0.52 & $<10^{-100}$ \\
 & Exposición en $v-1$ & $+0.72$ & 0.11 & $<10^{-11}$ \\
Rezagada: día anterior de Pleno & Exposición del día anterior & $-0.10$ & 0.30 & $0.75$ \\
\bottomrule
\end{tabular}
\end{center}
\vspace{0.15cm}
\begin{itemize}
  \item La mitad del efecto crudo era mecánica ($11.2$ vs $6.0$ del mundo barajado); \emphr{lo que sobra es enorme}, viaja más fuerte \emph{entre} comisiones, y deja un eco corto que muere a escala de días.
  \item \nota{Marco honesto (Manski, Angrist, Shalizi-Thomas): clustering de defección por líneas de red más allá de lo mecánico; influencia es la lectura natural, homofilia fina la alternativa.}
\end{itemize}
\end{frame}

\begin{frame}{¿Y quién transmite? Los novatos}
\begin{center}\small
\begin{tabular}{lccc}
\toprule
Canal del emisor & $\phi$ vecinos \emph{con} el atributo & $\phi$ \emph{sin} & dif.\ ($p$) \\
\midrule
Abogado & $+5.8$ & $+6.0$ & $0.88$ \\
Posgrado & $+4.3$ & $+7.0$ & $0.003$ \\
Experiencia institucional & $+2.6$ & $+9.3$ & $<0.001$ \\
\bottomrule
\end{tabular}
\end{center}
\vspace{0.25cm}
\begin{itemize}
  \item Contra la hipótesis de los emisores de élite: los abogados transmiten \emphr{igual} que el resto; los senior transmiten \emphr{menos}.
  \item La co-defección viaja por los \emphr{pares novatos}: la mayoría nueva se mueve junta; la vieja guardia no arrastra a nadie.
\end{itemize}
\end{frame}

\sectionslide{RQ3 — Textos}{¿Qué hace que un artículo sobreviva?}

\begin{frame}{RQ3: el modelo de supervivencia}
\begin{itemize}
  \item Unidad: el \emphr{artículo} (1.565, con coalición firmante $S_a$ de 2--16; 389 coaliciones). ¿Llega al borrador? (20.2\% lo logra.)
\end{itemize}
$$\Pr(sobrevive_a = 1) = \Lambda\Big(\alpha_c + \gamma_1\, \big|\bar\theta_{1,S_a} - \theta_{1,(103)}\big| + \gamma_2\, sd(\theta_{1,S_a}) + \gamma_3\, |S_a| + \boldsymbol{\gamma_4}^\top C_{S_a} + \boldsymbol{\gamma_5}^\top H_{S_a}\Big)$$
\begin{itemize}
  \item $C_{S_a}$ = bloque de red de la coalición (betweenness media, constraint media, densidad interna); $H_{S_a}$ = capital humano (prop.\ abogados, experiencia, grado medio); $\alpha_c$ = FE de comisión; EE por coalición; predictores estandarizados.
\end{itemize}
\end{frame}

\begin{frame}{RQ3: la tabla completa (tres modelos anidados)}
\begin{center}\footnotesize
\begin{tabular}{lccc}
\toprule
 & (1) Pivotal & (2) + Red & (3) + Capital humano \\
\midrule
\multicolumn{4}{l}{\emphr{--- Geometría ideológica de la coalición}}\\
Distancia de la coalición al pívot & $-0.35$ ($p=.001$) & $-0.48$ ($p=.009$) & $-0.61$ ($p=.002$) \\
Heterogeneidad ideológica $sd(\theta_1)$ & $+0.27$ ($p=.008$) & $+0.46$ ($p=.001$) & $+0.53$ ($p<10^{-4}$) \\
Tamaño de la coalición & $+0.20$ ($p=.12$) & $+0.18$ ($p=.12$) & $+0.15$ ($p=.21$) \\
\multicolumn{4}{l}{\emphr{--- Red de la coalición}}\\
Betweenness media & & $-0.22$ ($p=.083$) & $-0.23$ ($p=.076$) \\
Constraint media & & $+0.01$ ($p=.97$) & $+0.26$ ($p=.28$) \\
Densidad interna (pares con historia) & & $+0.25$ ($p=.024$) & $+0.32$ ($p=.002$) \\
\multicolumn{4}{l}{\emphr{--- Capital humano de la coalición}}\\
Prop.\ abogados & & & $-0.09$ ($p=.58$) \\
Prop.\ con experiencia previa & & & $-0.16$ ($p=.44$) \\
Grado académico medio & & & $-0.06$ ($p=.67$) \\
\midrule
AIC & 1393 & 1384 & 1386 \\
\bottomrule
\end{tabular}
\end{center}
\vspace{0.1cm}
\nota{Ganan las coaliciones bien paradas, anchas y con historia interna — no las tituladas. La densidad interna: de las parejas posibles entre los firmantes, ¿qué fracción ya había co-firmado en \emph{otras} iniciativas?}
\end{frame}

\begin{frame}{El premio a ensancharse, quintil a quintil}
\begin{center}\small
\begin{tabular}{lccc}
\toprule
Posición media de la coalición & Supervivencia & $sd(\theta_1)$ & $p$ \\
\midrule
Q1: izquierda $[-0.83, -0.73]$ & 7\% & $+1.18$ & $.050$ \\
\emphr{Q2: izquierda moderada} $[-0.73, -0.67]$ & 14\% & \emphr{$+2.46$} & $.001$ \\
Q3: centro-izquierda $[-0.66, -0.50]$ & 34\% & $+0.36$ & $.49$ \\
Q4: centro $[-0.50, -0.18]$ & 30\% & $-0.09$ & $.82$ \\
Q5: centro y derecha $[-0.17, +0.81]$ & 15\% & $+0.18$ & $.72$ \\
\bottomrule
\end{tabular}
\end{center}
\vspace{0.25cm}
\begin{itemize}
  \item El premio vive donde ensancharse \emphr{alcanza el pívot}: máximo en Q2 (desde $-0.7$ el brazo llega al 103); en el extremo, menor (ni estirándose); desde Q3, innecesario.
  \item Decaimiento monótono hacia cero: \emphr{la lógica de los 103 votos, por dentro}.
\end{itemize}
\end{frame}

% ================================================================ CIERRE
\sectionslide{Cierre}{La aritmética del 103}

\begin{frame}{La geometría era destino}
\fig{0.86}{survival_by_position.pdf}
\vspace{0.05cm}
\nota{Bandas: terciles de los 154 convencionales — como $103 = \tfrac{2}{3} \times 154$, el corte T2/T3 \emph{es} el pívot. El punto dulce no está sobre el pívot: está a su izquierda — el centro de masa de la mayoría que redactaba, con el brazo estirado hasta el voto 103.}
\end{frame}

\begin{frame}{La misma curva, persona a persona}
\fig{0.86}{success_by_ideology.pdf}
\vspace{0.05cm}
\nota{El éxito individual replica la geometría: máximo justo a la izquierda del pívot (listas locales, INN y Apruebo parados ahí); niveles moderados en la izquierda; piso plano en la derecha.}
\end{frame}

\begin{frame}{Síntesis}
\begin{itemize}
  \item En una asamblea de extraños, la colaboración se organizó con lo que la gente \emphr{traía puesto}: lista, ideología, territorio — no credenciales; y la estructura completa muestra coaliciones que \emphr{mezclan} género, educación y generación.
  \vspace{0.2cm}
  \item Ese capital \emphr{no movió posiciones} (selección, con el nulo defendido), pero \emphr{sí coordinó conducta} (la defección viaja por los lazos, transportada por los novatos) y \emphr{sí ganó textos} (equipos con historia, cerca del pívot, anchos).
  \vspace{0.2cm}
  \item Y la institución cierra el círculo: la regla de 2/3 fijó un pívot, y la red fue el mecanismo por el que esa aritmética se volvió comportamiento.
  \vspace{0.2cm}
  \item Lección de método: \emphr{la unidad de análisis real} (persona $\times$ documento) no es un tecnicismo — cambia respuestas.
  \vspace{0.3cm}
  \item[] \begin{center}\emphr{Se colaboró por cercanía; se ganó por amplitud.}\end{center}
\end{itemize}
\end{frame}

\begin{frame}{Limitaciones y trabajo en curso}
\begin{itemize}
  \item Las indicaciones aún no entran como red dinámica: el panel artículo $\times$ onda con "equipo de custodia" (firmantes génesis $\cup$ enmendadores) está diseñado, no corrido.
  \vspace{0.25cm}
  \item ERGM bipartito: EE por bootstrap de iniciativas y términos estructurales recién incorporados; sensibilidad de los decaimientos \emph{gw} pendiente.
\end{itemize}
\end{frame}

\begin{frame}[plain]
  \vfill
  \begin{center}
    {\Large ¡Gracias!}\\[2pt]
    {\color{ndred}\rule{0.30\textwidth}{1.1pt}}\\[0.7cm]
    \begin{tabular}{c@{\hspace{2.6cm}}c}
      \qrcode[height=2.4cm]{https://github.com/aoliveram/Networked-Influence-Chilean-Constitutional-Convention} &
      \qrcode[height=2.4cm]{https://github.com/aoliveram/constitutional-proposal-tracking} \\[6pt]
      {\small Proyecto de investigación (código y reporte)} & {\small Datos: constitutional-proposal-tracking} \\
    \end{tabular}\\[0.7cm]
    {\texttt{anibal.olivera.m@gmail.com}}
  \end{center}
  \vfill
\end{frame}

% ================================================================ BACKUP
\sectionslide{Backup}{Material de respaldo}

\begin{frame}{Backup: el modelo de eventos (RHEM) — la dinámica}
\begin{itemize}
  \item Las 947 iniciativas ocurren \emphr{en orden}. ¿Firmo contigo porque \emph{nos parecemos} o porque \emph{ya firmamos juntos}?
  \item Cada iniciativa: la coalición real contra 50 ficticias del mismo tamaño ese día. Estadística clave:
\end{itemize}
$$sub.rep^{(2)}(t,h) = \binom{|h|}{2}^{-1} \sum_{\{i,j\} \subseteq h} \#\{\text{eventos previos que contienen a } i \text{ y } j\}$$
\begin{center}\small
\begin{tabular}{lcc}
\toprule
 & Coef. & $p$ \\
\midrule
\emphr{Familiaridad de pares} ($sub.rep^{(2)}$) & $+2.11$ & $10^{-20}$ \\
Prop.\ pares misma lista & $+0.92$ & $10^{-8}$ \\
Dispersión ideológica $\theta_1$ & $-2.66$ & $10^{-39}$ \\
\bottomrule
\end{tabular}
\end{center}
\vspace{0.1cm}
\nota{La Convención se tejió \emphr{de a dos}: se reclutan duplas consolidadas. Y la homofilia sobrevive a la historia: parecerse y conocerse operan a la vez.}
\end{frame}

\begin{frame}{Backup: robustez del nulo de influencia (ventanas $\times$ dosis)}
\begin{itemize}
  \item 3 ventanas (completa / temprana $\leq$ mar / tardía $\geq$ abr) $\times$ 7 definiciones de exposición (acumulada, últimas $k$ ondas, decaimientos): $\hat\lambda \in [-0.01, +0.011]$; los pocos significativos de la tardía vienen con falsificación igual de significativa.
  \item Chequeo exploratorio con red imputada al primer mes (la ventana de los ``novatos susceptibles''): nada.
  \item Poder: la distancia real persona-vecindario (0.59) era espacio suficiente; MDE $= 0.011$.
\end{itemize}
\end{frame}

\begin{frame}{Backup: el modelo espacial del éxito (SDM) y sus impactos exactos}
\begin{itemize}
  \item Éxito acoplado en la red: $\rho = 0.89$ (Durbin, robusto 0.63--0.95 entre variantes de $W$).
  \item Impactos exactos con $(I-\rho W)^{-1}$ denso: directos chicos ($\approx$ OLS); totales amplificados $\sim 9\times$ por el multiplicador — la lectura de equilibrio literal no es creíble; $\rho$ se lee como composición + selección.
  \item Contrafactual tangible: pasar del vecindario P10 al P90 de éxito asocia $+0.078$ = \emphr{1.3 sd del éxito}.
\end{itemize}
\end{frame}

\begin{frame}{Backup: ¿por qué no un solo ERGM de las 947? (Simpson)}
\begin{itemize}
  \item Con intercepto único, las diferencias de composición \emph{entre} comisiones contaminan la homofilia \emph{dentro} de cada una: el agregado \emphr{invierte los signos} (misma lista pasa de $+0.2^{***}$ por comisión a $-0.09$ agregado).
  \item El modelo general correcto es la suite por comisión (o su versión conjunta con interceptos por red).
  \item MCMC de la espec extendida: infactible (semanas); MPLE + bootstrap de iniciativas: la ruta práctica.
\end{itemize}
\end{frame}

\begin{frame}{Backup: proporción de abogados por iniciativa (los junta el tema)}
\fig{0.9}{lawyer_share_initiatives.pdf}
\end{frame}

\begin{frame}{Backup: dinámica de retención textual}
\fig{0.9}{retention_dynamics_all_commissions.pdf}
\end{frame}

\begin{frame}{Backup: distribución de firmas}
\fig{0.9}{signature_distributions.pdf}
\end{frame}

\end{document}
